\documentclass[leqno]{article}
% Shared mark scheme preamble for resources.danieljsmith.org: the pack
% preamble, unchanged, then the scheme's own layer. A scheme's .tex is
% rendered from its YAML (scripts/README.md, Mark schemes), never edited.
\input{\CurrentFilePath/pack_preamble.tex}

% The page: equal side margins, so the table runs the full width of the
% header rule. The footers' right edge follows the text's.
\geometry{left=0.8in, right=0.8in}
\makeatletter
\renewcommand{\djs@footR}[1]{%
  \put(\dimexpr1in+\hoffset+\oddsidemargin+\textwidth\relax,
    -\dimexpr\paperheight-0.32in\relax){\makebox[0pt][r]{#1}}}
\makeatother

\usepackage[export]{adjustbox}
\usepackage{tabularray}
\definecolor{msGrid}{RGB}{176, 172, 164}
\definecolor{msTint}{RGB}{245, 243, 239}

% The header on every page, the front page included, and the booklet label
% under the contents-page title. \djsQuestionPageStyle stays empty.
\newcommand{\djsSchemeHeader}[1]{%
  \djsPackHeader[markscheme]{\textit{Mark scheme} - #1}%
  \renewcommand{\djsPackTitle}{#1}%
  \renewcommand{\djsBookletLabel}{Mark scheme}%
}

% The front page: the contents heading, spaced as tocloft spaces the packs'
% "Questions", then the contents table. A question's name and page number
% both link to its first page.
\newcommand{\msContentsHeading}{%
  \vskip\cftbeforetoctitleskip
  {\noindent\cfttoctitlefont\contentsname\par}%
  \vskip\cftaftertoctitleskip}
\NewTblrEnviron{mscontents}
\SetTblrInner[mscontents]{
  colspec = {Q[wd=26mm, l] X[l] Q[wd=16mm, c] Q[wd=14mm, r]},
  width = \linewidth,
  rowsep = 4pt, colsep = 6pt,
  hlines = {0.4pt, msGrid}, vlines = {0.4pt, msGrid},
  hline{2} = {2.2pt, scarlet},
  row{1} = {font=\small\bfseries},
  row{Z} = {font=\bfseries},
  column{2} = {font=\small},
}
\newcommand{\msContentsQuestion}[1]{\hyperref[q:#1]{Question #1}}
\newcommand{\msContentsPage}[1]{\hyperref[q:#1]{\pageref*{q:#1}}}
\newcommand{\msSep}{\hspace{0.9em}}

% The conventions pages' layout: headings in the documents' own bold, tables
% ruled and padded as the scheme's, roomy lists.
\newcommand{\msHead}[1]{%
  \par\vspace{18pt}%
  {\noindent\bfseries #1\par}%
  \vspace{6pt}}
\NewTblrEnviron{mstblr}
\SetTblrInner[mstblr]{rowsep = 6pt, colsep = 6pt,
  hlines = {0.4pt, msGrid}, vlines = {0.4pt, msGrid}}
\newenvironment{msList}
  {\begin{itemize}[leftmargin=1.2em, itemsep=5pt, topsep=3pt, parsep=0pt]}
  {\end{itemize}}

% A question's first page: the link target and the PDF bookmark, as
% \questionitem gives the packs.
\newcommand{\msQnum}{}
\newcommand{\msQuestion}[1]{%
  \clearpage
  \renewcommand{\msQnum}{#1}%
  \phantomsection\label{q:#1}%
  \addcontentsline{toc}{subsection}{\protect\djsTocEntry{#1}}}

% The scheme's table: Question 16mm, Scheme the rest, Marks 16mm, Notes
% 44mm; every cell ruled in msGrid; the scarlet rule under the head row. A
% long table has no caption or "continued" foot; a continuation page opens
% with a muted line above the repeated head row. The rules across a part's
% rows, the alternative's tint and the total row are the rendered table's.
\NewTblrTheme{ms}{%
  \DefTblrTemplate{firsthead}{default}{}%
  \DefTblrTemplate{middlehead,lasthead}{default}{%
    {\small\itshape\color{djsMuted}Question \msQnum\ continued\par}%
    \vspace{6pt}}%
  \DefTblrTemplate{firstfoot,middlefoot,lastfoot}{default}{}%
}
\NewTblrEnviron{msscheme}
\SetTblrOuter[msscheme]{long, theme=ms, entry=none, label=none}
\SetTblrInner[msscheme]{
  colspec = {Q[wd=16mm, l] X[l] Q[wd=16mm, c] Q[wd=44mm, l]},
  width = \linewidth,
  rowsep = 6pt, colsep = 6pt,
  rowhead = 1,
  vlines = {0.4pt, msGrid},
  hline{1} = {0.4pt, msGrid}, hline{2} = {2.2pt, scarlet},
  hline{Z} = {0.4pt, msGrid},
  column{1} = {font=\bfseries}, column{3} = {font=\bfseries},
  column{4} = {font=\small},
  row{1} = {rowsep=4pt, font=\small\bfseries},
}

% Two marks on a row: each note starts with its code.
\newcommand{\msNote}[2]{\textbf{#1}\enspace #2}
% A full-row figure: the worked solution's own picture, scaled uniformly and
% only down, to the table less the Question column. The picture is drawn
% with the line width it has in the solutions document, at its list depth
% there (the optional argument, 1 by default: a part's working), so one
% sized by \linewidth keeps its size; tabularray sets the cell's \linewidth
% to the cell's measured width, which would shrink it at every pass. The
% widths are measured once, in a solution box at the packs' text width
% (pack_preamble.tex's margins, 0.8in and 1in).
\newlength{\msSpanW}
\makeatletter
\newcommand{\ms@keepLW}[1]{%
  \expandafter\xdef\csname ms@solnLW@#1\endcsname{\the\linewidth}}
\AtBeginDocument{%
  \setlength{\msSpanW}{\dimexpr\textwidth-16mm-28pt\relax}%
  \setbox\z@\vbox{%
    \begin{minipage}{\dimexpr\paperwidth-1.8in\relax}
      \begin{djsSolution}
        \ms@keepLW{0}%
        \begin{enumerate}\item \ms@keepLW{1}%
        \begin{enumerate}\item \ms@keepLW{2}%
        \begin{enumerate}\item \ms@keepLW{3}%
        \end{enumerate}\end{enumerate}\end{enumerate}
      \end{djsSolution}
    \end{minipage}}}
\newcommand{\msFigure}[2][1]{%
  \adjustbox{max width=\msSpanW, max totalheight=0.6\textheight, valign=t}{%
    \hsize=\csname ms@solnLW@#1\endcsname
    \linewidth=\hsize
    #2}}
\makeatother
% A figure mark's feature, in the Scheme column.
\newcommand{\msFeature}[1]{\textit{#1}}
% An alternative's title row.
\newcommand{\msAlternative}[1]{\textit{Alternative:} #1}
% Drafts only: a queried mark's code, the query's why under its note, and a
% flag in its part's place.
\newcommand{\msQueried}[2]{#1~\(\to\) #2?}
\newcommand{\msQuery}[1]{{\itshape\color{djsMuted}Query: #1}}
\newcommand{\msFlag}[2]{\textbf{#1:}\enspace #2}


\djsSchemeHeader{Complex Series}
\djsSitePath{Core Pure}{Complex Numbers}

\djsLastUpdated{24 October 2026}
\begin{document}
\djsSetQuestionCount{7}
\djsAddMarks{1}{8}
\djsAddMarks{2}{8}
\djsAddMarks{3}{8}
\djsAddMarks{4}{9}
\djsAddMarks{5}{7}
\djsAddMarks{6}{12}
\djsAddMarks{7}{12}
\vspace*{6pt}
\djsContentsTitle
\vspace{30pt}
\msContentsHeading
\begin{mscontents}{}
Question & Marks by part & Marks & Page \\
\msContentsQuestion{1} & (a)\,4\msSep (b)\,4 & 8 & \msContentsPage{1} \\
\msContentsQuestion{2} & (a)\,4\msSep (b)\,4 & 8 & \msContentsPage{2} \\
\msContentsQuestion{3} & & 8 & \msContentsPage{3} \\
\msContentsQuestion{4} & (i)\,2\msSep (ii)\,7 & 9 & \msContentsPage{4} \\
\msContentsQuestion{5} & & 7 & \msContentsPage{5} \\
\msContentsQuestion{6} & (i)\,3\msSep (ii)\,9 & 12 & \msContentsPage{6} \\
\msContentsQuestion{7} & (i)\,3\msSep (ii)\,9 & 12 & \msContentsPage{7} \\
\SetCell[c=2]{r} Total & & 64 & \\
\end{mscontents}
\newpage
% The conventions pages, shared by every pack's mark scheme: one source of
% truth, \input after the front page, with no strand or pack logic. The keys
% on the first side, the rules on the second. The scheme layer defines
% \msHead, msList and mstblr.
\vspace*{6pt}%
{\noindent\cfttoctitlefont Reading the mark scheme\par}%
\vspace{10pt}%
These rules apply throughout, unless a note says otherwise.

\msHead{Marks}
Every mark is worth one, so \textbf{M1 A1} on a row is two marks. Only marks
coded ft follow through.\par
\vspace{10pt}%
\begin{mstblr}{
  colspec = {Q[wd=15mm, l] X[l]},
  width = \linewidth,
  hline{2} = {2.2pt, scarlet},
  column{1} = {font=\bfseries},
  row{1} = {font=\small\bfseries, rowsep=4pt},
}
Code & Meaning \\
M1 & Method: a correct method, applied to this question. A slip doesn't cost it if the method is clear. A formula must be used, not just quoted. \\
dM1 & Dependent method: as M1, but only if the method mark before it was earned, unless its note names another. \\
A1 & Accuracy: a correct answer or step. It needs the method mark it follows. \\
A1ft & As A1, following through from an earlier answer that the note names. \\
A1~AG & The last mark of a part whose answer is printed in the question: complete working, with no errors. \\
B1 & A correct value, fact or statement. It needs no method mark. \\
B1ft & As B1, following through from an earlier answer that the note names. \\
E1 & A written explanation, interpretation or conclusion. Concluding a proof by induction needs every earlier mark. \\
\end{mstblr}

\msHead{Abbreviations}
ft and AG qualify a code in the Marks column. The others appear in the
notes.\par
\vspace{10pt}%
\begin{mstblr}{
  colspec = {Q[wd=17mm, l] X[l] Q[wd=15mm, l] X[l]},
  width = \linewidth,
  column{1} = {font=\bfseries}, column{3} = {font=\bfseries},
}
ft & follow through & AG & answer given in the question \\
cao & correct answer only, with no follow-through & cso & correct solution only: no errors in the part \\
oe & or an equivalent form & awrt & any answer that rounds to this value \\
soi & seen, or implied by later working & isw & ignore working after a correct answer \\
SC & special case: credit outside the main method & their & the student's own earlier value \\
condone & a slip that doesn't cost the mark & [\;\,] & working the mark doesn't need \\
s.f. & significant figures & d.p. & decimal places \\
\end{mstblr}

\newpage
\msHead{Reading the scheme}
\begin{msList}
\item Each row is a line of working, with the marks it earns beside it.
  Where a row carries two marks, each note starts with its code.
\item \textbf{[\textit{n}]} ends each part with its total.
\item A shaded \textit{Alternative} is a second method for the same total.
  Mark an answer by the method it follows.
\item Any other valid method earns the equivalent marks, unless the question
  names a method or says \textit{hence}.
\item A sketch earns one B1 for each feature named below its figure.
\end{msList}

\msHead{Accuracy and answers}
\begin{msList}
\item Unless a question says otherwise, give answers to 3 significant
  figures, and angles in degrees to 1 decimal place. Angles in radians take
  3 significant figures.
\item A final answer with more figures is accepted if it rounds to the
  scheme's value; one with fewer is not. An answer spoiled by rounding too
  early loses its accuracy mark.
\item Where a question asks for an exact answer, a decimal loses the final
  accuracy mark. Otherwise either is accepted.
\item Wrong or missing units cost nothing unless the question asks for them.
\item A correct answer with no working earns full marks, unless the question
  needs the working: show that, prove, hence, a method it names, or an exact
  answer. A wrong answer with no working earns nothing. Calculator answers
  are treated the same way.
\item A correct answer, once reached, isn't spoiled by later working (isw).
\end{msList}

\msHead{Value of $\boldsymbol{g}$}
\begin{msList}
\item Take $g = 9.8\,\mathrm{m\,s^{-2}}$ unless a question states a value,
  which is then the only one accepted. Answers that use $g = 9.8$ may be
  given to 2 or 3 significant figures.
\item Where a question doesn't state $g$, an answer worked with $g = 10$
  throughout earns the same marks, keeping any value printed in the
  question. The worked solutions list the final answers with $g = 10$.
\end{msList}


\msQuestion{1}
\begin{msscheme}{
  hline{3-6,8-11} = {2-Z}{0.4pt, msGrid},
  hline{7,12} = {0.4pt, msGrid},
  row{12} = {font=\bfseries},
}
Question & Scheme & Marks & Notes \\
1(a) & {\(C + iS = e^{2i\theta} - \frac13 e^{5i\theta} + \frac19 e^{8i\theta} - \cdots\)} & M1 & {Uses \(\cos k\theta + i\sin k\theta = e^{ik\theta}\) to write at least three terms as exponentials.} \\
& {Geometric, first term \(e^{2i\theta}\), \(r = -\frac13 e^{3i\theta}\)} & A1 & {Correct common ratio, oe, soi by a correct sum.} \\
& {\(C + iS = \dfrac{e^{2i\theta}}{1 - \left(\their{-\frac13 e^{3i\theta}}\right)}\)} & M1 & {Uses \(\frac{a}{1-r}\) with their first term and their ratio of the form \(ke^{3i\theta}\).} \\
& {\(\phantom{C + iS} = \dfrac{3e^{2i\theta}}{3 + e^{3i\theta}}\)} & A1~AG & {Printed result from a correct sum without error. No check of \(|r| < 1\) needed.} \\
\nopagebreak & & [4] & \\
1(b) & {\(C + iS = \dfrac{3e^{2i\theta}(3 + e^{-3i\theta})}{(3 + e^{3i\theta})(3 + e^{-3i\theta})}\)} & M1 & {Multiplies numerator and denominator by the conjugate of the denominator, in exponential or trigonometric form.} \\
& {Denominator \(= 9 + 3\left(e^{3i\theta} + e^{-3i\theta}\right) + 1\)\\[4pt] \(\phantom{\text{Denominator}} = 10 + 6\cos 3\theta\)} & A1 & \\
& {Numerator \(= 9e^{2i\theta} + 3e^{-i\theta}\)\\[4pt] \(S = \dfrac{\mathrm{Im}\left(9e^{2i\theta} + 3e^{-i\theta}\right)}{10 + 6\cos 3\theta} = \ldots\)} & M1 & {Expands the numerator and takes its imaginary part over their real denominator. Condone sign slips.} \\
& {\(S = \dfrac{9\sin 2\theta - 3\sin\theta}{10 + 6\cos 3\theta}\)} & A1~AG & {Printed result from a correct numerator and denominator without error.} \\
\nopagebreak & & [4] & \\
\nopagebreak \SetCell[c=2]{r} Total for Question 1 & & 8 & \\
\end{msscheme}

\msQuestion{2}
\begin{msscheme}{
  hline{3-6,8-11} = {2-Z}{0.4pt, msGrid},
  hline{7,12} = {0.4pt, msGrid},
  row{12} = {font=\bfseries},
}
Question & Scheme & Marks & Notes \\
2(a) & {\(C + iS = \sum_{n=0}^{\infty}\left(-\tfrac12\right)^n e^{i(4n+1)\theta}\)} & M1 & {Combines the two series term by term, using \(\cos x + i\sin x = e^{ix}\), into one series of exponentials.} \\
& {Geometric, first term \(e^{i\theta}\), ratio \(-\tfrac12 e^{4i\theta}\)} & A1 & {Common ratio \(-\frac12 e^{4i\theta}\) oe, stated or used.} \\
& {\(C + iS = \dfrac{e^{i\theta}}{1 + \frac12 e^{4i\theta}}\)} & M1 & {Uses \(\frac{a}{1-r}\) with their first term and ratio, each of the form \(ke^{mi\theta}\).} \\
& {\(\phantom{C + iS} = \dfrac{2e^{i\theta}}{2 + e^{4i\theta}}\)} & A1~AG & {Correct sum to infinity leading without error to the printed result. \(|r| < 1\) need not be shown.} \\
\nopagebreak & & [4] & \\
2(b) & {\(C + iS = \dfrac{2e^{i\theta}(2 + e^{-4i\theta})}{(2 + e^{4i\theta})(2 + e^{-4i\theta})}\)} & M1 & {Multiplies numerator and denominator by \(2 + e^{-4i\theta}\) oe, such as \(2 + \cos 4\theta - i\sin 4\theta\).} \\
& {\((2 + e^{4i\theta})(2 + e^{-4i\theta}) = 5 + 4\cos 4\theta\)} & A1 & {Correct real denominator from a seen expansion.} \\
& {\(C = \mathrm{Re}\left(4e^{i\theta} + 2e^{-3i\theta}\right) \div \their{5 + 4\cos 4\theta}\)\\[4pt] \(\phantom{C} = \ldots\)} & M1 & {Expands the numerator to two exponential terms and takes the real part of each over their real denominator.} \\
& {\(C = \dfrac{4\cos\theta + 2\cos(-3\theta)}{5 + 4\cos 4\theta}\)\\[4pt] \(\phantom{C} = \dfrac{4\cos\theta + 2\cos 3\theta}{5 + 4\cos 4\theta}\)} & A1~AG & {Correct real parts over \(5 + 4\cos 4\theta\) leading without error to the printed result; \(\cos(-3\theta) = \cos 3\theta\) may be implied.} \\
\nopagebreak & & [4] & \\
\nopagebreak \SetCell[c=2]{r} Total for Question 2 & & 8 & \\
\end{msscheme}

\msQuestion{3}
\begin{msscheme}{
  hline{3-7} = {2-Z}{0.4pt, msGrid},
  hline{8} = {0.4pt, msGrid},
  row{8} = {font=\bfseries},
}
Question & Scheme & Marks & Notes \\
3 & {\(C + iS = \sum_{n=0}^{\infty} r^n e^{i(\alpha + n\beta)}\)\\[4pt] \(\phantom{C + iS} = e^{i\alpha}\sum_{n=0}^{\infty} \left(re^{i\beta}\right)^n\)} & M1 & {Writes the terms as \(r^n e^{i(\alpha + n\beta)}\) and identifies a geometric series with ratio \(re^{i\beta}\), soi by its use.} \\
& {\(C + iS = \frac{e^{i\alpha}}{1 - re^{i\beta}}\)} & M1 A1 & {\msNote{M1}{Uses \(\frac{a}{1 - q}\) with their first term and a complex ratio of the form \(re^{\pm i\beta}\). Convergence need not be stated.}\\ \msNote{A1}{oe}} \\
& {\(C + iS = \frac{e^{i\alpha}\left(1 - re^{-i\beta}\right)}{\left(1 - re^{i\beta}\right)\left(1 - re^{-i\beta}\right)}\)} & M1 & {Multiplies numerator and denominator by the conjugate of their denominator, in exponential or trigonometric form.} \\
& {\(\left(1 - re^{i\beta}\right)\left(1 - re^{-i\beta}\right)\)\\[4pt] \(\phantom{(1)} = 1 - r\left(e^{i\beta} + e^{-i\beta}\right) + r^2\)\\[4pt] \(\phantom{(1)} = 1 + r^2 - 2r\cos\beta\)} & A1 & \\
& {Equates real and imaginary parts of the numerator:\\[4pt] \(e^{i\alpha} - re^{i(\alpha - \beta)}\)\\[4pt] \(\phantom{e} = \cos\alpha - r\cos(\alpha - \beta)\)\\[4pt] \(\phantom{e = } + i\left(\sin\alpha - r\sin(\alpha - \beta)\right)\)} & dM1 & {Writes their numerator as real part plus \(i\) times imaginary part in terms of \(\alpha\) and \(\beta\), and equates these to \(C\) and \(S\). Condone sign slips.} \\
& {\(S = \frac{\sin\alpha - r\sin(\alpha - \beta)}{1 + r^2 - 2r\cos\beta}\)\\[4pt] \(C = \frac{\cos\alpha - r\cos(\alpha - \beta)}{1 + r^2 - 2r\cos\beta}\)} & A1 A1~AG & {\msNote{A1}{oe}\\ \msNote{A1~AG}{From the real part of a correct \(C + iS\) over the real denominator; working back from the printed result earns nothing.}} \\
\nopagebreak \SetCell[c=2]{r} Total for Question 3 & & 8 & \\
\end{msscheme}

\msQuestion{4}
\begin{msscheme}{
  hline{3-4,6-10} = {2-Z}{0.4pt, msGrid},
  hline{5,11} = {0.4pt, msGrid},
  row{11} = {font=\bfseries},
}
Question & Scheme & Marks & Notes \\
4(i) & {\(\sec\theta(\cos\theta + i\sin\theta)\)\\[4pt] \(\phantom{A} = \sec\theta\cos\theta + i\sec\theta\sin\theta\)} & M1 & {Expands the bracket, or factorises \(1 + i\tan\theta\), using \(\sec\theta = \frac{1}{\cos\theta}\) in at least one term.} \\
& {\(\phantom{A} = 1 + i\tan\theta\)} & A1~AG & {Both terms simplified without error to the printed result, from either side.} \\
\nopagebreak & & [2] & \\
4(ii) & {\(C + iS = 1 + \binom{n}{1}(i\tan\theta) + \binom{n}{2}(i\tan\theta)^2\)\\[4pt] \(\phantom{C + iS = 1} + \binom{n}{3}(i\tan\theta)^3 + \cdots\)} & M1 & {Combines the series as a binomial expansion in \(i\tan\theta\), showing at least the terms in \(\tan^2\theta\) and \(\tan^3\theta\) with their powers of \(i\).} \\
& {\(C + iS = (1 + i\tan\theta)^n\)} & A1 & \\
& {\(C + iS = \left[\sec\theta(\cos\theta + i\sin\theta)\right]^n\)\\[4pt] \(\phantom{C + iS} = \sec^n\theta(\cos\theta + i\sin\theta)^n\)} & M1 & {Substitutes part (i) into their \((1 + i\tan\theta)^n\) and takes the power of each factor.} \\
& {\(C + iS = \sec^n\theta(\cos n\theta + i\sin n\theta)\)} & dM1 A1 & {\msNote{dM1}{Applies de Moivre's theorem to \((\cos\theta + i\sin\theta)^n\).}} \\
& {\(S = \sec^n\theta\sin n\theta\)\\[4pt] \(C = \sec^n\theta\cos n\theta\)} & A1 A1~AG & {\msNote{A1}{\(S = \sec^n\theta\sin n\theta\) stated; the equating of imaginary parts may be implied.}\\ \msNote{A1~AG}{Real parts equated from a correct \(C + iS\), with no errors in the part.}} \\
\nopagebreak & & [7] & \\
\nopagebreak \SetCell[c=2]{r} Total for Question 4 & & 9 & \\
\end{msscheme}

\msQuestion{5}
\begin{msscheme}{
  hline{3-8} = {2-Z}{0.4pt, msGrid},
  hline{9} = {0.4pt, msGrid},
  row{9} = {font=\bfseries},
}
Question & Scheme & Marks & Notes \\
5 & {\(\sum_{r=0}^{n-1}\exp\left(2i\theta + \frac{4\pi ir}{n}\right)\)\\[4pt] \(\phantom{\sum} = \exp(2i\theta)\sum_{r=0}^{n-1}\left(\exp\left(\frac{4\pi i}{n}\right)\right)^r\)} & M1 & {Writes the sum as a constant times a sum of powers of one complex exponential. Condone slips in the ratio.} \\
& {Ratio \(\exp\left(\frac{4\pi i}{n}\right) \neq 1\) since \(0 < \frac{2}{n} < 1\)} & A1 & {Correct ratio, with a reason it is not 1 that uses \(n > 2\).} \\
& {\(\exp(2i\theta)\,\frac{1 - \left(\exp\left(\frac{4\pi i}{n}\right)\right)^n}{1 - \exp\left(\frac{4\pi i}{n}\right)}\)} & M1 & {Applies the geometric series sum formula with their first term and ratio. Condone a slip in the number of terms.} \\
& {\(\left(\exp\left(\frac{4\pi i}{n}\right)\right)^n = \exp(4\pi i) = 1 \Rightarrow\) sum \(= 0\)\\[4pt] \(\Rightarrow \sum_{r=0}^{n-1}\cos\left(2\theta + \frac{4\pi r}{n}\right) = 0\)\\[4pt] \(\phantom{\Rightarrow} \sum_{r=0}^{n-1}\sin\left(2\theta + \frac{4\pi r}{n}\right) = 0\)} & A1 & {Shows the ratio to the power \(n\) is 1, so the sum is 0, and that both its real and imaginary parts are 0.} \\
& {\(C = \sum_{r=0}^{n-1}\frac{1 + \cos\left(2\theta + \frac{4\pi r}{n}\right)}{2}\)\\[4pt] \(S = \frac12\sum_{r=0}^{n-1}\sin\left(2\theta + \frac{4\pi r}{n}\right)\)} & M1 & {Uses a double-angle identity to write \(C\) or \(S\) in terms of the real or imaginary part of the sum. Condone coefficient slips.} \\
& {\(C = \frac12(n) + \frac12(0) = \frac{n}{2}\)} & A1 & {cso, with \(\sum_{r=0}^{n-1} 1 = n\) seen or clearly used.} \\
& {\(S = \frac12(0) = 0\)} & A1~AG & {Correct identity applied to the zero imaginary part, reaching \(S = 0\) with no errors seen.} \\
\nopagebreak \SetCell[c=2]{r} Total for Question 5 & & 7 & \\
\end{msscheme}

\msQuestion{6}
\begin{msscheme}{
  hline{3-5,7-14} = {2-Z}{0.4pt, msGrid},
  hline{6,15} = {0.4pt, msGrid},
  row{15} = {font=\bfseries},
}
Question & Scheme & Marks & Notes \\
6(i) & {\((1 - 2e^{2i\theta})(1 - 2e^{-2i\theta})\)\\[4pt] \(\phantom{A} = 1 - 2e^{2i\theta} - 2e^{-2i\theta} + 4\)} & M1 & {Multiplies out to four terms, using \(e^{2i\theta}e^{-2i\theta} = 1\). Condone one sign slip.} \\
& {\(\phantom{A} = 5 - 2(e^{2i\theta} + e^{-2i\theta})\)} & A1 & {oe, such as \(5 - 2e^{2i\theta} - 2e^{-2i\theta}\).} \\
& {\(\phantom{A} = 5 - 2(2\cos 2\theta) = 5 - 4\cos 2\theta\)} & A1~AG & {Uses \(e^{2i\theta} + e^{-2i\theta} = 2\cos 2\theta\), or the cosine and sine forms, reaching the printed result without error.} \\
\nopagebreak & & [3] & \\
6(ii) & {\(C + iS = e^{i\theta} + 2e^{3i\theta} + \cdots + 2^{n-1}e^{(2n-1)i\theta}\)} & M1 & {Forms \(C + iS\) and writes its terms as exponentials; the first two terms and the last are enough.} \\
& {Geometric: first term \(e^{i\theta}\), ratio \(2e^{2i\theta}\)} & A1 & {Ratio \(2e^{2i\theta}\), soi; first term \(1\) accepted with \(e^{i\theta}\) factored out.} \\
& {\(C + iS = \frac{e^{i\theta}\left(1 - (2e^{2i\theta})^n\right)}{1 - 2e^{2i\theta}}\)} & M1 A1 & {\msNote{M1}{Uses \(\frac{a(1-r^n)}{1-r}\) or \(\frac{a(r^n-1)}{r-1}\) with their first term and ratio, of the form \(ke^{2i\theta}\), and \(n\) terms.}\\ \msNote{A1}{oe, such as numerator \(e^{i\theta} - 2^n e^{(2n+1)i\theta}\).}} \\
& {\(\phantom{C + iS} = \frac{e^{i\theta}(1 - 2^n e^{2ni\theta})(1 - 2e^{-2i\theta})}{(1 - 2e^{2i\theta})(1 - 2e^{-2i\theta})}\)} & M1 & {Multiplies numerator and denominator by the complex conjugate of their denominator.} \\
& {\(\phantom{C + iS} = \frac{e^{i\theta}(1 - 2^n e^{2ni\theta})(1 - 2e^{-2i\theta})}{5 - 4\cos 2\theta}\)} & A1 & {Denominator \(5 - 4\cos 2\theta\), from part (i); the numerator need not be expanded.} \\
& {Numerator \(= e^{i\theta} - 2e^{-i\theta} - 2^n e^{(2n+1)i\theta}\)\\[4pt] \(\phantom{\text{Numerator} =} + 2^{n+1}e^{(2n-1)i\theta}\)\\[4pt] Equates real and imaginary parts} & dM1 & {Expands their numerator to four terms, writes them in cosine and sine form and takes real and imaginary parts. Condone sign slips.} \\
& {\(S = \frac{2^{n+1}\sin((2n-1)\theta) - 2^n\sin((2n+1)\theta) + 3\sin\theta}{5 - 4\cos 2\theta}\)} & A1 & {oe.} \\
& {\(C = \frac{2^{n+1}\cos((2n-1)\theta) - 2^n\cos((2n+1)\theta) - \cos\theta}{5 - 4\cos 2\theta}\)} & A1~AG & {Real part of a correct expanded numerator, reaching the printed result without error.} \\
\nopagebreak & & [9] & \\
\nopagebreak \SetCell[c=2]{r} Total for Question 6 & & 12 & \\
\end{msscheme}

\msQuestion{7}
\begin{msscheme}{
  hline{3-5,7-14} = {2-Z}{0.4pt, msGrid},
  hline{6,15} = {0.4pt, msGrid},
  row{15} = {font=\bfseries},
}
Question & Scheme & Marks & Notes \\
7(i) & {\(1 - e^{2i\theta} = 1 - \cos 2\theta - i\sin 2\theta\)} & M1 & {Writes \(e^{2i\theta}\) in Cartesian form, as \(\cos 2\theta + i\sin 2\theta\) or by expanding \((\cos\theta + i\sin\theta)^2\).} \\
& {\(\phantom{1 - e^{2i\theta}} = 2\sin^2\theta - 2i\sin\theta\cos\theta\)} & M1 & {Uses \(1 - \cos 2\theta = 2\sin^2\theta\) and \(\sin 2\theta = 2\sin\theta\cos\theta\), or equivalent identities, on both parts. Condone sign slips.} \\
& {\(\phantom{1 - e^{2i\theta}} = 2\sin\theta(\sin\theta - i\cos\theta)\)} & A1~AG & {Printed result reached without error from correct identities, including the sign of the imaginary part.} \\
\nopagebreak & & [3] & \\
7(ii) & {\(C + iS = 1 - \binom{n}{1}e^{2i\theta} + \binom{n}{2}e^{4i\theta} - \cdots\)\\[4pt] \(\phantom{C + iS} = \sum_{r=0}^{n}\binom{n}{r}\left(-e^{2i\theta}\right)^r\)} & M1 & {Combines the series into terms \(\binom{n}{r}(-1)^r e^{2ri\theta}\), at least three seen or in general, and identifies a binomial expansion. Condone sign slips.} \\
& {\(C + iS = \left(1 - e^{2i\theta}\right)^n\)} & A1 & \\
& {\(C + iS = \left[2\sin\theta(\sin\theta - i\cos\theta)\right]^n\)\\[4pt] \(\phantom{C + iS} = 2^n\sin^n\theta\,(\sin\theta - i\cos\theta)^n\)} & M1 & {Substitutes the result of (i) into \(\left(1 - e^{2i\theta}\right)^n\) and separates the real factor \(2^n\sin^n\theta\).} \\
& {\(\sin\theta - i\cos\theta = -i(\cos\theta + i\sin\theta)\)} & M1 & {Writes \(\sin\theta - i\cos\theta\) as a multiple of \(\cos\theta + i\sin\theta\), or in exponential or modulus-argument form.} \\
& {\(\sin\theta - i\cos\theta = e^{-i\pi/2}(\cos\theta + i\sin\theta)\)} & A1 & {oe, such as \(e^{i(\theta - \pi/2)}\) or \(\cos\left(\theta - \frac{\pi}{2}\right) + i\sin\left(\theta - \frac{\pi}{2}\right)\).} \\
& {\((\sin\theta - i\cos\theta)^n = e^{-in\pi/2}\,e^{in\theta}\)} & dM1 & {Applies de Moivre's theorem, or the laws of indices, to raise their exponential or modulus-argument form to the power \(n\).} \\
& {\(C + iS = 2^n\sin^n\theta\, e^{i(n\theta - n\pi/2)}\)\\[4pt] \(\phantom{C + iS} = 2^n\sin^n\theta\bigl(\cos\left(n\theta - \tfrac{n\pi}{2}\right)\)\\[4pt] \(\phantom{C + iS = 2^n\sin^n\theta\bigl(} + i\sin\left(n\theta - \tfrac{n\pi}{2}\right)\bigr)\)} & A1 & {In exponential or modulus-argument form.} \\
& {\(S = 2^n\sin^n\theta\sin\left(n\theta - \frac{n\pi}{2}\right)\)\\[4pt] \(C = 2^n\sin^n\theta\cos\left(n\theta - \frac{n\pi}{2}\right)\)} & A1 A1~AG & {\msNote{A1}{oe.}\\ \msNote{A1~AG}{Equates real parts of a correct expression for \(C + iS\), reaching the printed form.}} \\
\nopagebreak & & [9] & \\
\nopagebreak \SetCell[c=2]{r} Total for Question 7 & & 12 & \\
\end{msscheme}

\end{document}
